\section{人机交互}

\subsection{互动与沟通}
\label{subsec:kommunikation}

互动与沟通这两个词经常混用，但严格来说不完全一样。

\emph{沟通（Kommunikation）} 指的是信息的传递或交换，可以是一个动作加一个反应，也可能只是单向的信息发出。而且沟通不一定是故意的，很多时候我们是无意识地在传递信息（比如你皱了下眉头，别人就接收到了）。

\emph{互动（Interaktion）}强调的是人与人之间的相互作用。互动一定是双向的：你做了动作，对方也做出反应，这一来一回构成了互动。

两者有很大的重叠，但互动更强调"来回"，沟通更强调"信息本身"。

\paragraph{Watzlawick 的五条沟通公理}

这是沟通学里最著名的理论之一，由心理学家 Paul Watzlawick 提出。
\begin{enumerate}
    \item 人不可能不沟通，哪怕你一言不发，你的沉默、表情、姿态本身也在传递信息。
    \item 每条信息都有两层含义：一是"内容层面"（你说的事实是什么），二是"关系层面"（你和听者之间是什么关系，语气亲密还是疏远）。
    \item 人类沟通同时用两种语言：一种是数字化的，逻辑清晰、有语法规则（比如说出来的话本身）；另一种是模拟化的，靠语气、表情、肢体动作传达关系层面的含义。
    \item 沟通是循环的因果关系——你说的话是对方反应的原因，对方的反应又成了你下一句话的原因，很难说谁是"起点"。
    \item 沟通要么是对称的，要么是互补的。对称是双方地位平等（比如朋友间的争论），互补是双方地位不对等（比如老师和学生）。
\end{enumerate}

\paragraph{Schulz von Thun 的沟通四层模型}

德国心理学家 Schulz von Thun 提出：你说的每一句话，其实同时包含四层信息：
\begin{enumerate}
    \item 事实信息：你客观陈述的内容是什么
    \item 自我表露：这句话无意中透露了你自己的什么状态或想法
    \item 关系暗示：这句话透露出你怎么看待对方、你们是什么关系
    \item 诉求：你说这句话，其实是想让对方做点什么
\end{enumerate}
经典例子是副驾驶对司机说前面是绿灯。表面上只是在陈述一个事实，但背后可能同时是我有点着急（自我表露）、我觉得你反应有点慢（关系暗示，可能让人不爽）、你可以走了（诉求）。同一句话，说话人的意图和听话人的理解经常对不上，这就是很多沟通误会的根源。

\paragraph{Jones \& Gerard 的四种社会互动模式}

这个理论描述了两个人互动时，各自到底在"跟随"什么：
\begin{enumerate}
    \item 完全互动：双方的目标和行为高度一致，配合得很好。
    \item 非对称互动：一方完全按自己的节奏来，不太理会对方。
    \item 伪互动：表面在互动，实际上双方都只顾着实现自己的目标，没有真正配合。
    \item 反应式互动：只是被动地对对方的刺激做出反应，没有自己的主动计划。
\end{enumerate}

这一节看似在讲人与人之间的沟通理论,但其实是为后面几节做铺垫。接下来会讲人和机器之间怎么互动，机器要模拟或应对人类沟通里的这些层次和模式，才能设计出好用的人机交互系统。理解人与人之间沟通的复杂性，是理解人机交互设计难点的第一步。

\subsection{人与机器的沟通}
\label{subsec:contactmanmachine}

沟通要发生，需要两个基本前提：
\begin{itemize}
    \item 共同的符号库（gemeinsamer Zeichenvorrat）：双方得用同一套语言或符号系统,不然收到信号也看不懂。比如只听得懂一种语言的人，收到用另一种语言的信息，是完全无法理解的。
    \item 通道和兼容的收发器官：得有一条信息能传过去的路径，双方的发送和接收设备也得匹配。
\end{itemize}
信号（Signal） 是信息的载体，可以是一个或多个符号，它的作用是在接收方那里触发某种反应。但同一个信号，不同接收者的反应可能完全不同，这取决于对方的知识、经验、教养、文化背景。人、动物、机器面对同一个信号，做出的反应逻辑也各不相同：人可能是本能反应，也可能是后天学来的行为模式和约定俗成的规则；机器则是靠预先设定的定义和规则来"反应"。

\paragraph{Shannon \& Weaver 的发送者-接收者模型}

这是信息论里最经典的模型之一,把一次沟通拆成几个环节：

信息源 → 发送者（编码） → 信号 → 传输通道 → 接收者（解码） → 目标

发送者把要传达的信息编码成信号，信号通过通信信道传输（传输过程通常会受到干扰），接收者把收到的信号解码还原成信息。

\emph{干扰（Störungen） 的来源}：
\begin{itemize}
    \item 信道里的信号叠加（噪声）
    \item 解码过程中的错误
    \item 双方符号库不完全兼容
\end{itemize}

\vspace{2em}

沿用 Shannon \& Weaver 的框架，\emph{沟通可以从三个层面来分析}，每个层面回答一个不同的问题：
\begin{itemize}
    \item \emph{技术层面（technical）}：符号有没有被正确传输？关注编码、解码和物理层面的信号传递本身。比如一个信号灯，是把电流转换成光信号、光波传播出去、再被视网膜转换成神经信号的过程。
    \item \emph{语义层面（semantic）}：传输的符号有没有精确传达出想表达的含义？关注符号的形式含义和内容组合。比如信号灯的颜色约定：红色=故障（加工中断）、黄色=警告（料仓存料少）、绿色=机器运行、蓝色=机器停止。
    \item \emph{效果层面（effectiveness）}：收到的信息有没有有效地引导出预期的行为？关注信息实际产生的效果。同样是信号灯，黄灯让员工去补料（有效），但如果员工看到红灯却在到处找故障原因而找不到（无效），说明信息虽然传到了，却没达成预期效果。
\end{itemize}
这三层递进的关系很重要：技术上传对了，不代表语义上双方理解一致；语义上理解一致，也不代表最终真的引导出了想要的行为。

沟通有三种主要形式：
\begin{itemize}
    \item 语言沟通（Verbale Kommunikation）：说出来的话、字母、数字、词语、句子和句法结构，主要负责传递事实信息，也带有一定的权重和强调。
    \item 副语言沟通（Paralinguistische Kommunikation）：不是“说了什么”，而是“怎么说”：语调、语速、停顿、笑声、叹气，这些提供了如何解读信息的线索。
    \item 非语言沟通（Nonverbale Kommunikation）：身体的动作和反应，包括触觉、身体距离、姿态、眼神接触、表情手势、着装风格等，主要用来表达情绪。
\end{itemize}
非语言沟通的难点在于它常常含糊不清、依赖上下文：同样是一声大喊，在庆典活动上是欢呼，在紧急情况下却是求救。脱离了场景，机器很难准确判断该如何理解。

工程标准上，沟通模型同样可以用 \ref{subsec:industrialcommunication} 小节中提到的 OSI 七层参考模型来表示。其中1、2层偏重传输本身，3层往上更偏向应用相关。在实际通信中，A站和B站之间的每一层都通过约定好的协议互相对应通信。发送方从第7层往下逐层封装，经过物理层传输后，接收方再从第1层往上逐层解封装还原。每一层之间都有明确定义的接口，各层之间互不干扰，这样才能让不同厂商、不同系统的设备也能互相听懂对方。

\subsection{技术可能性与限制}
\label{subsec:restriktoren}

前两节讲了沟通的理论模型。这一节我们落到实处：人有哪些发送器和接收器、机器有哪些对应的技术手段、这些渠道各自会遇到什么干扰，以及人和机器该如何互相迁就对方来实现有效沟通。

人的感官既是接收器，也是发送器。人身上天生自带一整套传感器（Sensorik）和执行器（Aktorik）：感知方面有视觉、听觉、嗅觉、味觉、平衡觉，以及触觉。表达方面，不同的感知渠道对应不同的表达方式：视觉对应表情和手势,听觉对应发声,触觉对应触摸、动作和身体姿态。人本身就是一套多模态的收发系统，这也是后面讨论机器如何与人沟通的基础——机器要跟人沟通，得对上人的这些感知渠道。

三种主要的传输方式：
\begin{itemize}
    \item 光学传输：机器靠指示灯、数码管显示、显示器发送，人靠眼睛接收。人靠表情、手势发送；机器靠摄像头接收。
    \item 声学传输：机器靠运行噪音、语音播放、蜂鸣器信号发送，人靠耳朵接收。人靠嘴巴说话、拍手发送，机器靠麦克风接收。
    \item 触觉传输：机器靠盲文显示屏、移动发送，人靠皮肤接收。人靠手指、手、身体发送，机器靠按键、键盘、触摸屏接收。
\end{itemize}
三种渠道的共同点是：人这一侧的干扰往往来自生理或环境限制，机器这一侧的干扰更多来自设备本身的性能或故障。

干扰可以按"沟通三层面"来分类。回顾 \ref{subsec:contactmanmachine} 节提到的技术层面、语义层面、效果层面，这里我们把\emph{常见干扰}具体归到了三个层级：
\begin{itemize}
    \item 技术层面的干扰：\emph{背景噪声}、信号衰减、信号失真、信道不兼容
    \item 语义层面的干扰：\emph{语言和手势不同}、编码方式不同、\emph{符号对应关系不同}、\emph{字符集不同}
    \item 效果层面的干扰：理解和意图相反、双方对世界的认知模型不同、信息含糊不清、\emph{缺乏上下文}、接收方信息不足
\end{itemize}
这套分类的意义在于：出了沟通问题，可以据此定位问题出在哪一层。是信号根本没传过去（技术层面）、传过去了但双方"翻译"不一致（语义层面）、还是双方都理解对了但没达成预期效果（效果层面）。

由于人和机器目前还没法做到完全对等地互相理解，因此我们需要双向的适应策略。
\begin{itemize}
    \item 人向机器靠拢：因为人机之间共同的"符号库"有限，交互方式也是由机器开发者预先设定好的，通常也不支持非语言沟通（手势、表情）,所以人往往需要调整自己的说话和行为方式，比如使用简化的、基于关键词的语言。例如直接说"Alexa，割草坪"智能音箱听不懂，得按规定格式说"Alexa，请割草坪"才能被识别执行。
    \item 机器向人靠拢：借助自学习算法扩展机器能识别的符号库、用人工智能方法识别非语言信息，具体体现在把多个沟通渠道的信息语义化地组合起来分析、让机器学习原本没有预设的新反应、以及对自然语言进行理解和解释。这是让机器越来越"像人"、越来越好用的方向。
\end{itemize}

\paragraph{机器的非语言表达}

机器也要学会打手势或做表情等非语言表达方式。这是为了通过额外的沟通渠道减少误解，更好地向人传达机器的意图，建立人对机器的信任，以及让机器感淡化，交互更自然。

具体做法包括给机器配上模拟的脸，模仿人类的手势和表情，营造出接近人类行为的印象。这样一来，人会不自觉地用理解人的方式去理解机器的行为。

\paragraph{人机交互的四种模式}

我们用一个“真实-虚拟”的矩阵把人机交互分成四种模式，如表 \ref{tab:realvirtualmatrix} 所示。

\begin{table}[htbp]
    \centering
    \caption{真实-虚拟矩阵}
    \label{tab:realvirtualmatrix}
    \begin{tabular}{lcc}
        \toprule
        - & 真实 & 虚拟 \\
        \midrule
        人 & 对话系统 & Human-in-the-Loop \\
        机器 & Hardware-in-the-Loop & 仿真 \\
        \bottomrule
    \end{tabular}
\end{table}

人和真实机器之间的交流是对话系统；人进入一个虚拟环境里参与和影响系统运行,就是 Human in the Loop。真实的硬件参与到一个虚拟系统的运行验证中,是 Hardware in the Loop。而双方都在虚拟世界里,就是纯粹的仿真。

\subsection{应用示例}

前面几节讲了沟通理论、人和机器各自的收发渠道、以及会遇到的干扰。这一节用六个真实的工业应用案例，把这些理论串起来。每个案例都标明了用了哪种沟通渠道（光学/声学/触觉）、方向是单向还是双向，以及具体靠哪些身体部位和设备来实现。

\paragraph{触摸屏与机械按钮}

触摸屏和机械按钮用于设备和机器人控制系统的调试与操作。如图 \ref{fig:touchtaster} 所示。工作方式上，设备通电靠机械开关来完成初始启动；急停这类关键功能用机械按钮，保证可靠性；显示屏上则是可以灵活定制、随需求变化的功能界面；设备状态和运行情况可以直接观察到，更详细的状态信息和故障信息则通过显示屏和指示灯呈现。

\begin{figure}[htbp]
    \centering
    \begin{subfigure}[b]{0.47\textwidth}
        \centering
        \includegraphics[width=\textwidth]{touchtaster1.png}
        \caption{固定式触摸屏}
        \label{fig:touchtaster1}
    \end{subfigure}
    \hfill
    \begin{subfigure}[b]{0.48\textwidth}
        \centering
        \includegraphics[width=\textwidth]{touchtaster2.png}
        \caption{移动式触摸屏}
        \label{fig:touchtaster2}
    \end{subfigure}
    \caption{触摸屏与机械按钮}
    \label{fig:touchtaster}
\end{figure}

用到的沟通渠道：光学通道（机器→人，负责状态展示）和触觉通道（人→机器，负责操作输入）。具体载体是触摸屏、机械按钮和开关、指示灯（有些和按钮做在一起），人这边用的是眼睛和手。

这个例子说明了最基础的人机交互方式。精简、可靠，关键安全功能坚持用机械按钮而不是纯触屏，也体现出讲义前面提到的技术层面干扰要尽量降到最低的设计思路。

\paragraph{数字化工位辅助系统（GeSA）}

数字化工位辅助系统用在手工装配工位，给工人提供数字化的辅助引导。如图 \ref{fig:gesa} 所示。它能自动识别装配进度，工人可以通过手势和二维码卡片来操作辅助系统，产品的中间检验和最终检验也是通过光学方式完成的。

\begin{figure}[htbp]
    \centering
    \begin{subfigure}[b]{0.44\textwidth}
        \centering
        \includegraphics[width=\textwidth]{gesa1.png}
        \caption{GeSA的布置}
        \label{fig:gesa1}
    \end{subfigure}
    \hfill
    \begin{subfigure}[b]{0.49\textwidth}
        \centering
        \includegraphics[width=\textwidth]{gesa2.png}
        \caption{操作界面}
        \label{fig:gesa2}
    \end{subfigure}
    \caption{GeSA}
    \label{fig:gesa}
\end{figure}

用到的沟通渠道是双向的光学通道：机器一侧靠深度摄像头和彩色摄像头识别，人这边用眼睛看显示屏、用手做手势，摄像头也能看到工件本身。

相比触摸屏与机械按钮，这里机器已经能主动看懂人的手势动作，是一个更主动、更智能的交互方式，呼应了 \ref{subsec:restriktoren} 小节提到的机器向人靠拢的适应策略。

\paragraph{装配辅助系统（Mona）}

装配辅助系统，如图 \ref{fig:mona} 所示，是面向大尺寸产品装配过程中的工人引导。系统靠骨骼追踪（Skelettracking）技术识别工人手臂姿态和手部位置，工作信息显示在显示屏上，工人通过两段式的手臂手势来操作系统。

\begin{figure}[htbp]
  \centering
  \includegraphics[width=0.4\textwidth]{mona.png}
  \caption{Mona}
  \label{fig:mona}
\end{figure}

沟通渠道同样是双向光学通道，用到深度摄像头、彩色摄像头、显示屏，人这边靠眼睛观察、手臂做动作。

这个案例和 GeSA 很像，但交互的身体部位从手扩展到了手臂，说明针对不同的装配场景（大件产品需要更大幅度的动作），交互方式也需要相应调整。

\paragraph{协作机器人的语音控制}

协作机器人用于协作装配场景。机器人按工序一步步地拿取和翻转装配对象，工人则负责抓取和安装零件，如图 \ref{fig:sprachsteuerung} 所示。工作说明和交互选项显示在屏幕上，工人可以用语音指令操作辅助系统，系统还能通过图像识别装配进度并进行产品检验。

\begin{figure}[htbp]
    \centering
    \begin{subfigure}[b]{0.48\textwidth}
        \centering
        \includegraphics[width=\textwidth]{sprachsteuerung1.png}
        \caption{协作机器人的布置}
        \label{fig:sprachsteuerung1}
    \end{subfigure}
    \hfill
    \begin{subfigure}[b]{0.48\textwidth}
        \centering
        \includegraphics[width=\textwidth]{sprachsteuerung2.png}
        \caption{机器人视角}
        \label{fig:sprachsteuerung2}
    \end{subfigure}
    \caption{协作机器人}
    \label{fig:sprachsteuerung}
\end{figure}

这里首次出现了声学通道（人→机器，靠语音），同时叠加了双向光学通道（显示屏 + 彩色摄像头 + 深度摄像头，配合机器人动作展示状态）。用到的设备包括显示屏、机器人动作、麦克风阵列、深度摄像头,人这边用眼睛、声音。

这是六个案例里第一次用到说话这种更自然的沟通方式，也呼应 \ref{subsec:kommunikation} 小节提到的语言沟通。比起按钮和手势，语音交互对人来说更省心，但也对机器的语音识别能力提出了更高要求。

\paragraph{虚拟示教（KoKoMo）}

虚拟示教是对协作装配过程进行虚拟化呈现的系统。工作方式是先在虚拟环境中对机器人进行示教（Teach-in），再把学到的动作轨迹传输到真实设备上执行，生成的程序还能回传到仿真环境中验证。整个装配过程可以在 VR 和 AR 中先行测试。

\begin{figure}[htbp]
    \centering
    \begin{subfigure}[b]{0.56\textwidth}
        \centering
        \includegraphics[width=\textwidth]{kokomo1.png}
        \caption{硬件视角}
        \label{fig:kokomo1}
    \end{subfigure}
    \hfill
    \begin{subfigure}[b]{0.40\textwidth}
        \centering
        \includegraphics[width=\textwidth]{kokomo2.png}
        \caption{软件视角}
        \label{fig:kokomo2}
    \end{subfigure}
    \caption{虚拟示教（KoKoMo）}
    \label{fig:kokomo}
\end{figure}

这个案例用到了三种通道：双向光学通道（VR眼镜呈现画面）、双向触觉通道（VR控制器提供操作和反馈）、声学通道（耳机→人，输出提示音）。用到的设备有VR眼镜、VR控制器、耳机、红外传感器。

这是六个案例里交互渠道最丰富的一个，也正好对应 \ref{subsec:restriktoren} 小节提到的 Human in the Loop。人通过多种感官同时沉浸在虚拟环境里，跟真实设备的动作数据打通。

\paragraph{半机械化的机械手焊接}

工业焊接场景的应用：焊枪装在轻型机器人上，机器人架在带齿轮驱动的导轨上（相当于增加了第七轴以扩大作业范围），导轨通过磁铁固定在工件上；操作人员大致示教焊枪的轨迹点位和角度；车间里所有工人腰间都配有移动式急停按钮；通过专用摄像头和显示屏监控熔池状态；用操纵杆修正轨迹和焊枪角度。

沟通渠道是光学通道（机器→人，通过专用摄像头显示熔池状态）和触觉通道（人→机器，通过急停按钮和操纵杆操作）。人这边用眼睛看、手操作。

这个案例强调的是安全性：移动式急停按钮让车间里所有人都能随时中止流程，这也是工业场景下人机交互设计必须优先考虑的安全冗余。

\vspace{2em}

把六个案例放在一起看会发现几个明显的趋势：
\begin{itemize}
    \item 用到的沟通渠道从简单（案例一：触摸屏+按钮）逐步叠加更多通道（案例五：光学+触觉+声学三合一）
    \item 光学通道几乎是每个案例的标配，因为"看"是最直接、最容易实现双向反馈的方式
    \item 安全相关的功能（急停）无论系统多智能，始终倾向于用最可靠的机械式触觉交互来实现，而不是依赖摄像头识别或语音识别
    \item 案例越靠后，机器"理解"人的能力越强（从固定按钮，到识别手势，到语音理解，再到多通道融合的沉浸式交互），呼应第4节讲的"机器向人靠拢"的发展方向
\end{itemize}



